\documentclass[10pt]{article}
\usepackage[letterpaper,margin=0.6in]{geometry}
\usepackage[T1]{fontenc}
\usepackage{amsmath,amssymb}
\usepackage{booktabs}
\usepackage{enumitem}

% Compact sections without requiring titlesec.
\makeatletter
\renewcommand\section{\@startsection{section}{1}{\z@}%
  {-1.5ex \@plus -.4ex \@minus -.2ex}%
  {0.5ex \@plus .2ex}%
  {\normalfont\bfseries}}
\makeatother
\setlist[itemize]{leftmargin=1.1em,topsep=1pt,itemsep=1pt,parsep=0pt}
\setlength{\parskip}{2pt}
\pagestyle{empty}

\begin{document}
\small

\begin{center}
  {\large\bfseries Recreating Figure 1 of Schultz, Dayan \& Montague (1997)}\\[1pt]
  {\small Anthony Avila \quad\textbullet\quad \texttt{anthony\_avila@berkeley.edu}
  \quad\textbullet\quad code and figures in \texttt{td\_dopamine.ipynb}}
\end{center}

\section*{The claim and the model}

The paper argues that dopamine neurons do not code for the occurrence of a reward, but for an error
between the predicted and the actual time and magnitude of that reward. A monkey presses a lever
after a light cue (CS) and receives fruit juice (US). Before training its dopamine neurons burst when
the juice arrives; after training the burst moves to the cue; and if the juice is then withheld,
firing is depressed below baseline at exactly the moment it was due.

Reward is the immediate signal $r(t)$, while value is expected total future reward, estimated by
$\hat{V}(t)$. The quantity the paper identifies with dopamine output is the temporal-difference
error $\delta(t) = r(t) + \gamma\hat{V}(t+1) - \hat{V}(t)$, which is available at every step rather
than only at the end of a trial. A single weight per cue could not encode \emph{when} the reward is
due, so the cue is given a complete serial-compound representation ($x_i(t) = 1$ exactly $i$ steps
after cue onset and $0$ otherwise), with $\hat{V}(t) = \sum_i w_i x_i(t)$ and a per-trial update
$\Delta w_i = \alpha \sum_t x_i(t)\,\delta(t)$.

I run 80 steps per trial, cue at $t=10$, reward of size $1$ at $t=60$ (the times in the paper's Fig.~3), 50 delay features, $\alpha = 0.1$, $\gamma = 1$, and 650 training trials, across 8 seeds.
Four sources of non-determinism are added: juice volume ($\sigma = 0.15$), reward omitted on 10\% of
training trials, random weight initialisation ($\sigma = 0.01$), and noise on the reported error
($\sigma = 0.05$).

\section*{What I found}

All three effects appear. Predicted values assume perfect learning; measured values are mean
$\pm$ s.d.\ over the 8 seeds.

\begin{center}
\footnotesize
\begin{tabular}{@{}llcccc@{}}
\toprule
& & \multicolumn{2}{c}{$\delta$ at cue ($t=10$)} & \multicolumn{2}{c}{$\delta$ at reward ($t=60$)}\\
\cmidrule(lr){3-4}\cmidrule(lr){5-6}
Case & Intuition & predicted & measured & predicted & measured\\
\midrule
Naive, juice delivered   & ``Didn't see that coming'' & $0$  & $\approx 0$        & $+1$ & $+1.064 \pm 0.065$\\
Trained, juice delivered & ``Just as I thought''      & $+1$ & $+0.889 \pm 0.021$ & $0$  & $+0.118 \pm 0.047$\\
Trained, juice withheld  & ``Where's the juice?''     & $+1$ & unchanged          & $-1$ & $-0.914 \pm 0.103$\\
\bottomrule
\end{tabular}
\end{center}

\noindent
The learner stores \emph{when} the juice was due, not merely that some was coming: the omission trace
is flat from the cue all the way to $t=60$ and only then dips, with no intervening event to trigger
it. The cue response is unchanged in the last row because $\delta(10)$ depends only on the weights.
Spread across seeds is small relative to the effects: the largest s.d.\ is about a tenth of the
reward magnitude.

\section*{Where my version differs from the paper's, and why}

\begin{itemize}
  \item The predicted reward is not perfectly silent ($+0.118$, not $0$), and the cue burst
  undershoots at $+0.889$. Both follow from the 10\% omission rate during training: $\hat{V}$
  converges to the \emph{expected} reward, not to $1$. Setting that rate to zero gives $+0.995$ at
  the cue, $+0.017$ at the reward and $-1.007$ on omission, matching the paper's middle panel. It is
  a consequence of the non-determinism I chose, not an error in the learner.
  \item My naive panel is not the paper's naive panel. Fig.~1 (top) is labelled ``(No CS)'':
  no cue is presented at all and the histogram is aligned on the reward. Mine presents the cue with
  untrained weights, so it is the first trial of training rather than an unpredicted reward in
  isolation. The traces coincide only because the initial weights are near zero.
  \item One cue, not two. The caption of Fig.~3 gives cues at steps 10 \emph{and} 20 with
  reward at 60; I kept the first. Its body text calls the same figure a reward ``20 time steps into
  the future,'' disagreeing with its own caption; I followed the caption.
  \item $\gamma = 1$ is load-bearing. Across a 50-step interval, discounting erases the
  effect the figure exists to show: $\gamma = 0.99$ shrinks the cue burst to $+0.59$ and
  $\gamma = 0.95$ to $+0.08$. Fig.~3 reports weights that ``eventually all saturate to 1,'' implying
  $\gamma = 1$ there too.
  \item Errors are symmetric in $\delta$ but not in spikes. The dip, $-0.914$, nearly
  mirrors the naive burst, $+1.064$, yet a neuron cannot fire at a negative rate, so the depression
  in Fig.~1 is bounded below by the basal rate. My raster panel reproduces this by suppressing
  background firing, which is why the dip looks shallower than the burst.
\end{itemize}

\section*{What didn't work}

\begin{itemize}
  \item Varying the seed alone. A deterministic TD learner on a complete serial compound
  follows exactly one trajectory: with all four noise magnitudes at zero the 8 seeds return
  byte-identical arrays and the figure collapses to a single line. Showing spread required making the
  task and the report stochastic, not just reseeding.
  \item Discounting. $\gamma < 1$ is the textbook default and the wrong choice at this
  cue--reward interval, for the reason above.
  \item Training only as long as the stimulus is wide. The burst does not jump to the cue;
  it drifts backwards one feature at a time at a rate set by $\alpha$. With $\alpha = 0.1$ the peak
  of $\delta$ is still at $t=55$ after 50 trials and reaches the cue near trial 400, hence the 650
  trials. Training briefly leaves a burst stranded mid-interval, which looks like a bug and is not.
\end{itemize}

\section*{One thing I'd try next}

\looseness=-1
Recreate Fig.~3 rather than Fig.~1: plot $\delta$ as a surface over time step \emph{and} trial index,
so the burst is seen migrating from reward to cue instead of sampled at three snapshots; the
training loop already stores every trial's error. Sweeping $\alpha$ against the cue--reward interval
would then show whether that backward drift rate is a real prediction of the model or an artefact of
how the cue is represented.

\end{document}
